% SPDX-License-Identifier: Apache-2.0
\section{Conclusion}

An undocumented binary format was worked out from experiments alone,
to the point where a reader reproduces a declared ground truth for
1049 minimal cases and for six designs, the largest of which records
18\,875\,466 value changes over 5696 objects, and where a converter
produces FST files that a standard viewer opens.

Three parts of the method carried the result. Minimal pairs, so that a
difference between two databases has one cause. A ground truth written
from the design before the database is opened, so that a decoding
claim can fail. And guards outside the project, so that the checker
and the decoder cannot share a mistake.

Two directions remain open. A second Vivado version would separate
properties of the format from properties of one release, which is the
largest single gap in what can be claimed. A large SystemVerilog
design would do for that language what NEORV32 did for VHDL: the
corpus covers SystemVerilog only in miniature, and the two defects
found by large designs were both found in the parts of the format that
only a large design exercises.

\section*{How it is built, and why}

Everything here is built by Bazel: the Vivado installation, every
corpus simulation, the reader, the converter, the C library it calls,
and this report. That is not a preference for a build tool. The method
of Section~\ref{sec:method} rests on a case differing from another
case in one way, so the rest of the environment has to be pinned:
the simulator version, the source files, the compile order, the flags,
and the tool that reads the result. Bazel gives declared inputs,
cached outputs and nothing rebuilt without a reason, which is what
makes a difference between two databases a difference in the design
rather than in the machine that produced it.

The installation is provisioned into a shared cache by
\texttt{rules\_vivado}, so the simulator itself is a build input
rather than something a person installed. The corpus is 1049
simulation targets, and one command re-runs them all. The C library
of Section~\ref{sec:output} is built with a Bazel-fetched toolchain,
so the release binary does not depend on the compiler of the machine
that built it.

The reasoning behind that setup, and the pinning it requires, is
written up on the author's site: hermetic, ephemeral and reproducible
builds~\cite{her}, the whole hardware stack under one build
system~\cite{btd}, an ephemeral machine to run them
on~\cite{homelab}, and the Vivado rules
themselves~\cite{rulesvivado}.

\section*{Availability}

The repository, the corpus, the format documents and the tool are at
\url{https://git.hdlfactory.com/HDL/wdbcvt}. This report is built from
\texttt{//docs/latex} in that repository and is attached to each
release.

The forge is private, and the reason is the runner. A build here needs
a machine with the Vivado installation and the shared caches on it,
because every corpus case is a simulation; a hosted runner cannot do
that. The repository is therefore where the runner is.

If you would like the source in a public place and are willing to
sponsor a runner that can hold a Vivado installation, please get in
touch through \url{https://hdlfactory.com}. That is the one thing
standing between this work and a public mirror.
